\section{The actual local period rings and the bare Galois stage}
\label{sec:local-stages}

\status{This section reconstructs the local-ring and invariant calculations
needed by \texttt{LocStage}, and the actual empty-local-index construction
of \texttt{NStage}. Its period rings are the specified Witt, divided-power,
and adic constructions. The final stage is attached to the same base,
columns, and adapted row record. Its two index types are empty; it supplies
bare stage existence, without an arithmetic twist reading, a nonzero
projector, or a comparison theorem for an arbitrary geometric Galois
representation. These are human mathematical proofs from the pinned
source, without a Lean build or a kernel-verification claim.}

\paragraph{Conventions.}
Throughout LS1--LS17, $p$ is any rational prime, including $2$,
$\Omega=\overline{\mathbb Q}_p$, $C=\widehat\Omega$, $O=\mathcal O_C$,
and $G=\operatorname{Gal}(\Omega/\mathbb Q_p)$. We normalize the
absolute value by $|p|=p^{-1}$. For a finite extension
$K\subset\Omega$, put $G_K=\operatorname{Gal}(\Omega/K)$ and
$K_0=K\cap\mathbb Q_p^{\mathrm{ur}}$, where
$\mathbb Q_p^{\mathrm{ur}}$ is the field generated by roots of unity of
order prime to $p$. Completion, finite-extension valuation theory,
finite-field Galois theory, and the defining Witt identities are used in
their ordinary mathematical senses. NC1--NC2 already provide the
complete-DVR and normalized-valuation facts. All further period-ring
properties used below are proved here.

\paragraph{LS1: an explicit Ax approximation and completed fixed fields.}
If $H\le G$, $a\in\Omega$, and $|ha-a|\le r$ for every $h\in H$,
then there is $b\in\Omega^H$ such that
\[
                     |a-b|\le16r.                 \tag{LS.1}
\]
Here is the finite-orbit proof, including the uniform constant. If the
orbit has at most $2^k$ elements, let
\[
 C_k=2^{\,4-(2k+4)/2^k},\qquad
 d_k=2^{(k+1)/2^k};\quad C_0=1,\quad C_{k+1}=C_kd_k\le16.
 \tag{LS.2}
\]
Induct on $k$. The one-point case is immediate. Otherwise let $T$ be
the orbit, $n=|T|$, with $2^k<n\le2^{k+1}$, and form
$f(X)=\prod_{z\in T}(X-z)$. Its Hasse derivative
$P=\partial^{[n-2^k]}f$ is $H$-fixed, has degree $2^k$, and leading
coefficient $\binom n{n-2^k}$. Expansion at $a$ gives
\[
 |P(a)|\le r^{2^k},\qquad
 \left|\binom n{n-2^k}\right|\ge n^{-1}.          \tag{LS.3}
\]
The second estimate follows from Legendre's factorial formula: the
$p$-valuation of a binomial coefficient is a sum of zeros and ones,
with no terms above $\lfloor\log_p n\rfloor$.
Factor $P$ over $\Omega$. Some root $z$ satisfies
$|z-a|^{2^k}\le n r^{2^k}$, hence $|z-a|\le d_kr$.
Every $H$-conjugate of $z$ is a root of $P$ and
$|hz-z|\le d_kr$ by the ultrametric inequality. Induction applied to
$z$ gives an $H$-fixed $b$ with
$|z-b|\le C_kd_kr$. The same inequality bounds $|a-b|$.
Every algebraic element has a finite orbit, proving (LS.1).

Consequently
\[
                         C^H=\overline{\Omega^H}.   \tag{LS.4}
\]
The inclusion from right to left follows from isometry of the action.
For the other inclusion, approximate $x\in C^H$ by $a\in\Omega$
within $\delta$; then $|ha-a|\le\delta$ for every $h\in H$.
Equation (LS.1) gives a fixed algebraic point within $16\delta$ of
$x$. In particular $C^{G_K}=K$, since a finite extension of
$\mathbb Q_p$ is complete and therefore closed in $C$.

\paragraph{LS2: a uniform cyclotomic projection bound.}
Choose compatible primitive roots $\zeta_n$ of order $p^n$, with
$\zeta_0=1$, $\zeta_{n+1}^p=\zeta_n$, and put
$K_n=\mathbb Q_p(\zeta_n)$. For $n\ge1$, $m\ge0$, $N=p^m$, and
$z=\zeta_{n+m}$, every $y\in K_{n+m}$ has a unique expansion
\[
 y=\sum_{i=0}^{N-1}a_i z^i,\quad a_i\in K_n,
 \qquad |p|\,|a_0|\le |y|.                         \tag{LS.5}
\]
We give the bound without a trace-integrality assumption. For a
primitive $p^r$-th root $z$, all primitive conjugates have the same
distance from $1$: the geometric-sum identity gives the inequality
in each direction. Evaluating the cyclotomic product at $1$ yields
\[
              |z-1|^{e_r}=p^{-1},\qquad
              e_r=p^{r-1}(p-1).                    \tag{LS.6}
\]
For $Q\in\mathbb Q_p[X]$ of degree $<e_r$, the nonzero terms of
$Q(z-1)$ have distinct valuations modulo $\mathbb Z$. Thus its norm
is the largest term norm, and
$p^{-1}|Q_j|\le|Q(z-1)|$ for every coefficient. Taylor substitution
$X\mapsto X+1$ and its inverse have integral binomial coefficients;
therefore the same estimate holds for the coefficients of a
polynomial evaluated at $z$. This also proves that
$1,z,\ldots,z^{e_r-1}$ are independent, so the cyclotomic degree is
exactly $e_r$.

Write each $a_i=P_i(\zeta_n)$, with $\deg P_i<e_n$, and set
$P(X)=\sum_{i<N}P_i(X^N)X^i$. Its degree is $<Ne_n=e_{n+m}$.
Its coefficients at powers divisible by $N$ are exactly those of
$P_0$. Applying the preceding estimate and then the ultrametric
inequality to $P_0(\zeta_n)$ proves (LS.5).
The extension $K_{n+m}/K_n$ has degree $N$, and averaging its roots
shows that $a_0=N^{-1}\operatorname{Tr}_{K_{n+m}/K_n}(y)$.
These projections are compatible as $m$ varies, commute with $G$,
and have norm at most $p$, independently of $m$ and $n$.
They extend to maps
\[
 R_n:\widehat{\bigcup_{r\ge1}K_r}\longrightarrow K_n,
 \qquad R_nx\longrightarrow x.                    \tag{LS.7}
\]
Indeed they are eventually the identity on the dense union, and the
uniform bound extends this convergence to its completion. This proof
also covers $p=2$, $n=1$, where $K_1=\mathbb Q_2$.

\paragraph{LS3: every nonzero Tate twist has zero invariants.}
Let $\chi:G\to\mathbb Z_p^\times$ be the character determined by
$g\zeta_n=\zeta_n^{\chi(g)\bmod p^n}$. For any finite $K$ and any
nonzero integer $i$,
\[
 \bigl\{x\in C:g(x)\chi(g)^i=x\quad(g\in G_K)\bigr\}=0.
 \tag{LS.8}
\]
First take $K=\mathbb Q_p$. The kernel of $\chi$ fixes precisely
the cyclotomic union in $\Omega$, so LS1 puts such an $x$ in the
completion in (LS.7). Equivariance gives
$g(R_nx)\chi(g)^i=R_nx$. For each $n$ there is
$g\in G_{K_n}$ with $\chi(g)^i\ne1$. Otherwise the image of this
group under $\chi$ would consist of roots of $X^{|i|}-1$, hence be
finite, contradicting the unbounded degrees
$[K_{n+m}:K_n]=p^m$ proved in LS2. This $g$ fixes $R_nx\in K_n$,
so $R_nx=0$ for every $n$. Taking the limit proves the assertion.

For general $K$, let $d=[G:G_K]$ and set
$f(g)=g(x)\chi(g)^i$. The hypothesis makes $f$ constant on cosets
$gG_K$. Thus $Y=\prod_{gG_K}f(g)$ satisfies
$h(Y)\chi(h)^{id}=Y$ for every $h\in G$, by permutation of the
cosets. The already proved case gives $Y=0$. A factor is zero, and
all character factors are units, so $x=0$. No estimate involving
$[K:\mathbb Q_p]^{-1}$ is needed.

\paragraph{LS4: the actual Fontaine map and its principal kernel.}
Set
\[
 R=\varprojlim_{x\mapsto x^p}O/pO,\quad A=W(R),\quad
 \theta\left(\sum_{j\ge0}p^j[x_j]\right)
       =\sum_{j\ge0}p^j x_j^\sharp .                \tag{LS.9}
\]
The sharp map is multiplicative; its reduction modulo $p$ is the
zeroth inverse-limit coordinate. It is not asserted to be additive.
The Witt identities make (LS.9) a ring homomorphism. Every element $y\in O$ has a $p$-th root in $O$:
approximate $y$ by algebraic elements $y_n$ and successively
choose roots $x_n^p=y_n$ with
$|x_{n+1}-x_n|\le|y_{n+1}-y_n|^{1/p}$, using the factorization
of $X^p-y_{n+1}$ at $x_n$. The ultrametric inequality makes
$x_n$ Cauchy when $y_n$ is Cauchy, and its limit is the required
root. Iterating gives compatible $p$-power roots of $y$, hence
a preimage under sharp. Thus $\theta$ is onto. The ring $R$ is a perfect valuation domain:
its norm is $|x|=|x^\sharp|$, and if $|y|\le|x|$, the ratios of
the exact sharp-root coordinates of $y$ and $x$ give an element of
$R$ whose product with $x$ is $y$. A zero sharp value forces every
coordinate to be zero.

Choose $\varpi\in R$ with $\varpi^\sharp=p$. Then
\[
                 \ker\theta=(p-[\varpi]).          \tag{LS.10}
\]
To prove this, write $a=[a_0]+pb$. If $\theta a=0$, then
$|a_0|\le|p|$, so $a_0=\varpi c$ in the valuation ring $R$.
Hence
$a=(p-[\varpi])(-[c])+p([c]+b)$, and applying $\theta$ shows
$[c]+b$ is again in the kernel. Iterating and summing the
$p$-adically convergent coefficients proves divisibility. The
reverse inclusion is immediate. The nonzero zeroth Witt coefficient
of $p-[\varpi]$ also makes this element a nonzerodivisor.

For the compatible cyclotomic system of LS2 let
$\varepsilon\in R$ have sharp-root coordinates $\zeta_n$, and put
\[
 \mu=[\varepsilon]-1,\quad
 \nu=[\varepsilon^{1/p}]-1,\quad
 \xi=\sum_{j=0}^{p-1}[\varepsilon^{1/p}]^j,
 \qquad \mu=\xi\nu.                                \tag{LS.11}
\]
Then $\theta\xi=0$, and $\xi$ is a unit multiple of the generator
in (LS.10). To check the unit, reduce a factorization
$\xi=(p-[\varpi])u$ modulo $p$. The zeroth coefficient of $\xi$
is $(\varepsilon^{1/p}-1)^{p-1}$ and has norm $|p|$, so $|u_0|=1$;
a Witt vector with unit zeroth coefficient is a unit. For odd $p$,
the norm assertion follows from (LS.6), since the error between
sharp of a difference and the difference of sharp values is in
$pO$, of strictly smaller norm. For $p=2$ take one more inverse
Frobenius root: $\zeta_2^2=-1$, so
$|(\zeta_2-1)^2|=|2|$, whereas $|\zeta_2-1|>|2|$.
The same reduction argument followed by squaring gives
$|\varepsilon^{1/2}-1|=|2|$. This checks the dyadic case explicitly.

The quotient $A/(\xi^n)$ is $p$-torsion-free and $p$-adically
complete for every $n$. Indeed $\xi\bmod p\ne0$ in the domain
$R$, so $p^k\mid\xi^n b$ implies $p^k\mid b$. This makes
$(\xi^n)$ closed in the complete ring $A$: a Cauchy sequence of
multiples of $\xi^n$ has Cauchy quotients. Compatible quotient
classes modulo $p^k$ admit successively adjusted lifts with
successive differences in $p^kA$, proving completeness. These
arguments also show $p$-saturation of every $(\xi^n)$.

\paragraph{LS5: the actual de Rham DVR, filtration, and algebraic section.}
Put $A_P=A[1/p]$ and
\[
 B_{\mathrm{dR}}^+=\varprojlim_n A_P/(\xi^n),\qquad
 B_{\mathrm{dR}}=B_{\mathrm{dR}}^+[1/\xi],\qquad
 \operatorname{Fil}^aB_{\mathrm{dR}}=\xi^aB_{\mathrm{dR}}^+.
 \tag{LS.12}
\]
These agree with the source's completion and localization, since its
localized theta kernel is $(\xi)$, and every other generator differs
by a unit. The map $A_P\to B_{\mathrm{dR}}^+$ is injective.
In fact $\xi^n\mid a$ for every $n$ forces $a_0=0$ modulo $p$,
because $|\xi_0|<1$; after division by $p$, saturation permits the
same argument repeatedly, and $p$-adic separation gives $a=0$.
Clearing a power of $p$ proves the assertion on $A_P$.

The completion is separated and complete for $(\xi)$, with residue
field $C$, and multiplication by $\xi$ is injective. For the last
claim, lift an element modulo $\xi^{n+1}$ and cancel $\xi$ in
$A_P$; the resulting class modulo $\xi^n$ vanishes. An element of
the completion with nonzero residue is a unit, by the geometric
series lifting of its inverse. Every nonzero element is therefore
$\xi^r$ times a unit, for a unique $r\ge0$. This proves that
$B_{\mathrm{dR}}^+$ is a complete DVR and $B_{\mathrm{dR}}$ is its
fraction field. The filtration in (LS.12) is decreasing,
multiplicative, exhaustive and separated. The Galois action on
$C$, its valuation ring, the inverse-limit coordinates, the Witt
vectors, and then the completion preserves $\ker\theta$ and all
its powers. Thus it preserves every filtration step and commutes
with the residue map $\theta_D:B_{\mathrm{dR}}^+\to C$.

There is a unique $\mathbb Q_p$-algebra map
\[
 o:\Omega\longrightarrow B_{\mathrm{dR}}^+,
 \qquad\theta_D(o(a))=a,                            \tag{LS.13}
\]
where $\mathbb Q_p$ enters by the actual map $\mathbb Z_p\to A$.
For each algebraic $a$, lift the simple residue root $a$ of its
minimal polynomial by Newton iteration in the complete DVR.
The derivative has nonzero residue and is a unit, so the lift is
unique. On each finite separable subextension this is a field
embedding; uniqueness makes the embeddings compatible, additive,
multiplicative, and $G$-equivariant. In particular it agrees with
the prescribed constants from $\mathbb Z_p$. This construction does
not identify the residue map with a coefficient section on all of
$C$; the section asserted here has the algebraic domain $\Omega$.

\paragraph{LS6: the divided-power completion embeds in the actual de Rham ring.}
Inside $A_P$, set
\[
 D_0=A[\gamma_n:n\ge0],\qquad \gamma_n=\xi^n/n!,
 \qquad D=\widehat{D_0}^{\,(p)}.                    \tag{LS.14}
\]
Since $\ker\theta=(\xi)$, this is the source's divided-power
envelope $A_{\mathrm{cris},0}$, and $D=A_{\mathrm{cris}}$.
The multiplication identity
$\gamma_i\gamma_j=\binom{i+j}{i}\gamma_{i+j}$ shows that $D_0$
is the $A$-span of the displayed generators. It is $p$-torsion-free.
Its $p$-adic completion remains $p$-torsion-free, has dense image
of $D_0$, and satisfies
\[
 p^r\mid s\text{ in }D\quad\Longleftrightarrow\quad
 p^r\mid s\text{ in }D_0\qquad(s\in D_0).          \tag{LS.15}
\]
These assertions follow by projecting the inverse limit to
$D_0/p^rD_0$; to cancel $p$, lift one level farther and cancel in
the torsion-free ring $D_0$.

Modulo $\xi^nA_P$, the image $M_n$ of $D_0$ contains
$A/(\xi^n)$ and satisfies
\[
 p^n M_n\subset A/(\xi^n)\subset M_n.              \tag{LS.16}
\]
Indeed the terms $\gamma_j$ with $j\ge n$ disappear, while
$p^n/j!\in\mathbb Z_p$ for $j<n$. By LS4 the smaller module is
$p$-adically complete; the sandwich makes $M_n$ complete and
separated for the same topology. More explicitly, a Cauchy sequence
in $M_n$, after removal of a fixed initial term, has a tail in the
smaller module, and the two induced topologies differ by at most
$n$ powers of $p$. Extend $D_0\to M_n$ to its completion.
Uniqueness makes these maps compatible in $n$, giving
\[
 j:D\longrightarrow B_{\mathrm{dR}}^+,
 \quad j|_{D_0}\text{ is the actual inclusion},\quad
 p^n j(x)\bmod\xi^n\in\operatorname{im}(A).        \tag{LS.17}
\]
This is the literal boundedness field \texttt{IsCrisP.bdd}, with
$N=n$. Any map with the stated extension and boundedness is unique:
modulo $\xi^n$ its values lie in a fixed $p$-power enlargement of
$A/(\xi^n)$; dense approximations and separation there force equality.

We prove injectivity, rather than inferring it from density. Let
$L_N=\sum_{j<N}A\gamma_j$. The elementary relation
$\xi\gamma_j=(j+1)\gamma_{j+1}$ gives
\[
 L_N\cap\xi^N A_P\subset A\,(N\gamma_N).          \tag{LS.18}
\]
For induction, write $s=a\gamma_N+y$ with $y\in L_N$ and assume
$\xi^{N+1}\mid s$. The preceding induction step writes
$y=cN\gamma_N$. Multiplication by $N!$ and cancellation of
$\xi^N$ show that $\xi\mid a+Nc$ in $A_P$, hence in $A$ by
$p$-saturation. The displayed relation proves the next step.

Suppose $j(x)=0$, and write $x=s+pz$, $s\in D_0$.
For every $n$, (LS.17) and a sufficiently accurate dense approximation
give another expression $x=u+pz'$ with
$u\in D_0\cap\xi^nA_P$. Here is the adjustment explicitly:
write $x=s'+p^{n+1}v$ and choose $a\in A$ with
$p^n j(v)\equiv a\pmod{\xi^n}$; then $u=s'+pa$ has zero image
modulo $\xi^n$, and $x-u$ is divisible by $p$ in $D$.
Equation (LS.15) gives $s=u+pw$, $w\in D_0$.
Choose $M$ with $s\in L_M$, take $n=pM$, and split $w$ into
its terms of index $<n$ and its remaining terms. Thus
$s-pw_{<n}\in L_n\cap\xi^nA_P$. By (LS.18) this difference
is $a n\gamma_n$, divisible by $p$ in $D_0$. Hence $p\mid s$,
so $x=py$. As $p$ is a unit in $B_{\mathrm{dR}}^+$, also
$j(y)=0$. Iteration puts $x$ in every $p^rD$, whence $x=0$.

The actual Galois action preserves $D_0$ and extends to $D$.
The Witt Frobenius $\varphi$ does so as well: if
$\theta a=0$, then $a^p\in pD_0$ and
$\varphi(a)=a^p+pb\in pD_0$, so
$\varphi(a)^n/n!\in D_0$. The resulting actions commute, and
$j$ is Galois equivariant by its bounded uniqueness. No continuity
from the $p$-adic topology on $D$ to the $\xi$-adic topology on the
whole de Rham ring is presumed.

\paragraph{LS7: the logarithmic period, with integral convergence at $2$.}
For any $e\in R$ with $e^\sharp=1$, let $m=[e]-1$.
There is $w\in D_0$ such that
\[
                    pw=m^{p-1},\qquad w^p\in pD_0. \tag{LS.19}
\]
First suppose the first inverse-Frobenius sharp value is a nontrivial
$p$-th root. Put $X=[e^{1/p}]$ and $\eta=\sum_{i<p}X^i$.
The characteristic-$p$ polynomial identity gives
$m^{p-1}=\eta^p+pz$ with $z\in A$; applying $\theta$ shows
$\theta z=0$. Thus $w=\eta^p/p+z\in D_0$.
The divided powers of $\eta$ and $z$ show $w^p\in pD_0$:
the term $\eta^{p^2}/p^p$ is integral and divisible by $p$, since
$v_p((p^2)!)=p+1$, and the other terms follow from the binomial
coefficients and $z^p\in pD_0$. If the first nontrivial sharp root
occurs later, descend to it and use
$[e]-1=\xi(e)([e^{1/p}]-1)$ repeatedly. If there is no such root,
all coordinates of $e$ equal $1$, and (LS.19) holds with $w=0$.

Define $c_k=m^k/(k+1)$ in $A_P$. These elements lie in $D_0$.
For $q=\lfloor k/(p-1)\rfloor$, $r=k-q(p-1)$ and
$v=v_p(k+1)$, we have $v\le q$ and
\[
 c_k=(\text{a }p\text{-adic unit})\,p^{q-v}m^r w^q,
 \qquad p^{\lfloor q/p\rfloor}\mid c_k.            \tag{LS.20}
\]
The inequality $v(p-1)+1\le p^v\le k+1$ proves $v\le q$;
the divisibility follows from (LS.19). For $k>0$, also
$c_k\in\sqrt{pD_0}$. Therefore
\[
 u_e=\sum_{k\ge0}(-1)^k c_k\in D^\times,
 \qquad t_e=m u_e=\log[e]\in D.                   \tag{LS.21}
\]
Convergence follows from (LS.20). Modulo $p$, the difference
$u_e-1$ is a finite sum of nilpotents, hence nilpotent; lifting its
inverse through the $p$-adic completion proves that $u_e$ is a unit.
This argument uses no odd-prime exception. Multiplying a partial
sum by its factorial denominator gives exactly the finite sums in
the source's \texttt{IsLogA}, so (LS.21) has its stated meaning.

In $B_{\mathrm{dR}}^+$, $j(t_e)$ is the formal logarithm of
$1+j(m)$: modulo $\xi^n$, its terms of degree at least $n$
vanish, and (LS.17) transports the integral limiting relation.
Formal power-series identities and injectivity of $j$ give
$t_{e^a}=a t_e$ for integers $a$, and $\varphi(t_e)=pt_e$.
For the fixed primitive system $\varepsilon$, put $t=t_\varepsilon$.
It is nonzero and
\[
       \varphi(t)=pt,\qquad g(t)=\chi(g)t,
       j(t)\in\operatorname{Fil}^1\setminus\operatorname{Fil}^2.
 \tag{LS.22}
\]
For the Galois formula, choose integers $a_n\equiv\chi(g)\pmod{p^n}$.
The element $g\varepsilon\,\varepsilon^{-a_n}$ has its first
$n$ sharp-root coordinates equal to $1$; its $p^n$-th root $e_n$
therefore has sharp value $1$. Additivity of formal logarithms gives
$g(t)-a_nt=p^nt_{e_n}\in p^nD$. Completeness and separation give
the asserted limit. Finally $\mu=\xi\nu$ and
$\theta_D\nu=\zeta_1-1\ne0$; together with $u_\varepsilon$ a unit,
this proves the filtration assertion.

\paragraph{LS8: the actual crystalline core and unramified constants.}
For the chosen $\varepsilon$ take $B=D[1/\mu]$.
Since $t=\mu u_\varepsilon$, this is $D[1/t]$ as well;
the relation $pw=\mu^{p-1}$ makes $p$ a unit in $B$.
The injection $j$ extends uniquely to
\[
                    \iota:B\hookrightarrow B_{\mathrm{dR}}. \tag{LS.23}
\]
LS13 below checks the identification with all localization generators
in the source, so this notation does not substitute a formal field
model for its actual ring.

For a root $c\in\Omega$ of order $m$ prime to $p$, choose the unique
compatible $p$-roots within its finite group of $m$-th roots. They
define $c^\flat\in R$ with $(c^\flat)^m=1$ and
$(c^\flat)^\sharp=c$. In the complete DVR, both
$j([c^\flat])$ and $o(c)$ are roots of $X^m-1$ with residue $c$.
Their difference times the geometric sum is zero, and the latter
has nonzero residue $mc^{m-1}$. Thus
\[
                       j([c^\flat])=o(c).          \tag{LS.24}
\]
The subring of $B$ generated by $\mathbb Q_p$ and these lifts is a
field: every nonzero element is algebraic over $\mathbb Q_p$, so
its inverse is a polynomial in it, by its minimal polynomial with
nonzero constant term. Hence (LS.24) extends to a unique embedding
$u:\mathbb Q_p^{\mathrm{ur}}\to B$ with $\iota u=o$.

Reduction is injective on roots of unity of order prime to $p$.
The residue field of $\Omega$ is $\overline{\mathbb F}_p$; Hensel
lifting $X^{p^d}-X$ identifies each finite residue field with its
unramified extension. Its residue Frobenius therefore lifts to
compatible automorphisms, giving
$\sigma:\mathbb Q_p^{\mathrm{ur}}\to\mathbb Q_p^{\mathrm{ur}}$
with $\sigma(c)=c^p$ on these roots. Every finite intermediate
extension is cyclic unramified, so $\sigma$ preserves $K_0$ and
is bijective there. On its valuation ring it reduces to Frobenius.
Uniqueness on the generating roots gives
\[
 \varphi(u(c))=u(\sigma c),\qquad
 g(u(c))=u(gc),\qquad\iota(gb)=g(\iota b).         \tag{LS.25}
\]
The compatible roots, the tilt, $o$, $j$, $\iota$, $u$, $t$,
$\sigma$, and $\chi$ are now all on these same rings. The actions
extend to $B$, since (LS.22) makes the images of $t$ units there.
These give all maps and laws required by \texttt{Core}; LS13 below
identifies this localization with the source's precise period ring
and transports the construction. The Galois and Frobenius maps
are induced on coordinates and completions, and the coefficient
embeddings are the ones just constructed.

\paragraph{LS9: Frobenius is injective after divided-power completion.}
This needs a completion argument; injectivity on $A$ alone would
not suffice. Write $\xi=w+pa$ with $w=[\xi_0]$ a Teichm\"uller
element. Then $\varphi(w)=w^p$, $w\bmod p\ne0$, and
\[
 D_0\subset M=A[w/p]=A[\xi/p]\subset A_P,
 \qquad p^nM\bmod\xi^n\subset\operatorname{im}(A).
 \tag{LS.26}
\]
The first inclusion follows from $p^j/j!\in\mathbb Z_p$; the
last follows by expanding a polynomial in $\xi/p$ and separating
its terms of degree less than $n$.

We record the finite coefficient calculation used with this larger
ring. If $s\in M$ and, for every $r$, there are $u,v\in M$ with
\[
              \varphi(s+p^k u)=p^r v,
 \quad\text{then}\quad s\in p^kM.                 \tag{LS.27}
\]
Here $\varphi$ on $A_P$ is the extension of the bijective Witt
Frobenius. For completeness, put $I=(p,w)$. Regularity of $w$
modulo $p$ gives
$px\in I^{N+1}\Rightarrow x\in I^N$: expand
$I^{N+1}=Aw^{N+1}+pI^N$, reduce modulo $p$, then cancel $p$.
Every $v\in M$ has a denominator-clearing expression
$p^Nv=c$, $c\in I^N$.
For a polynomial $h=\sum h_iX^i$ define
\[
 F_0=0,\quad F_{N+1}=\varphi(h_N)w^{pN}+pF_N,
 \quad
 pF_N=p^N\varphi\!\left(\sum_{i<N}h_i(w/p)^i\right).
 \tag{LS.28}
\]
The precise induction used here is the following, with any integer
$l\ge1$ and recurrence $F_{N+1}=a_Nw^{lN}+pF_N$:
\[
 b w^{lN}+F_N\in I^{N+r},\quad i<N,\quad
 l(i+1)\le i+r+1
 \quad\Longrightarrow\quad a_i\in(p,w^l).             \tag{LS.28a}
\]
Induct on $N$. For the step from $N$ to $N+1$, put
$D=N+r+1$ and $A'=bw^l+a_N$. The hypothesis is
$A'w^{lN}+pF_N\in I^D$. If $D\le lN$, remove the first
term, which already belongs to $I^D$, and apply $p$-saturation
to obtain $F_N\in I^{N+r}$. The induction hypothesis with
$b=0$ handles $i<N$; the terminal $i=N$ is impossible since
its required inequality would give $l(N+1)\le D\le lN$.
If $D>lN$, put $e=D-lN>0$. Expand
$I^D=Aw^D+pI^{D-1}$ to write
\[
 A'w^{lN}+pF_N=w^Dc+pG,\qquad G\in I^{N+r}.
\]
Reduction modulo $p$ and cancellation of the regular element
$w^{lN}$ give $A'-w^ec=pd$ for some $d\in A$.
Cancel $p$ in the displayed equation to obtain
$d w^{lN}+F_N=G$. The induction hypothesis with $b=d$
handles $i<N$. At $i=N$, the required inequality is exactly
$e\ge l$, and
$a_N=pd+w^ec-bw^l\in(p,w^l)$. This proves (LS.28a)
with all its range conditions.

Now suppose $\varphi(h(w/p))=p^{p(i+1)+1}v$ with $v\in M$.
Write $p^{N_0}v=c\in I^{N_0}$ and take
$N=N_0+\deg(h)+i+1$ (the zero polynomial is immediate).
The last equality in (LS.28) gives
\[
 F_N=c\,p^{\deg(h)+i+1+p(i+1)}\in I^{N+p(i+1)}.
\]
Apply (LS.28a) with $l=p$, $r=p(i+1)$, $b=0$, and
$a_j=\varphi(h_j)$. Its inequalities hold, so
$\varphi(h_i)\in(p,w^p)$. Bijectivity of Witt Frobenius,
$\varphi(p)=p$, and $\varphi(w)=w^p$ then give
$h_i\in(p,w)$.

For $k=1$, apply this calculation to $h=f+p g$, where
$s=f(w/p)$ and $u=g(w/p)$ in the hypothesis of (LS.27), choosing
$r=p(i+1)+1$ separately for each coefficient $f_i$.
It follows that every $f_i\in(p,w)$, and substitution
$w=p(w/p)$ gives $s\in pM$. Repeated cancellation of $p$
proves (LS.27) for every $k$.

Now suppose $x\in D$ and $\varphi x=0$. Write
$x=s+p^kz$ with $s\in D_0$. Approximate $z$ arbitrarily closely
by elements of $D_0$. Applying Frobenius, and using (LS.15), gives
the hypothesis of (LS.27), so $s\in p^kM$.
For fixed $n$ choose $k=n+r$. By (LS.26) and (LS.17),
\[
 j(x)\bmod\xi^n
   =j(s)+p^{n+r}j(z)\bmod\xi^n
     \in p^r\bigl(A/(\xi^n)\bigr).
\]
This holds for every $r$; separation in LS4 makes the class zero.
Thus $j(x)=0$ and LS6 gives $x=0$. Frobenius on $D$ is injective.
It is injective on $B$ as well: clear a power of $t$, apply
Frobenius, and cancel the nonzero denominator using $D\hookrightarrow B$.

\paragraph{LS10: the integral Frobenius fixed ring is $\mathbb Z_p$.}
One has
\[
                         D^{\varphi=1}=\mathbb Z_p. \tag{LS.29}
\]
First $\varphi(D_0)\subset A+pD_0$. On a positive divided power
$a^n/n!$, $\theta a=0$, this follows from
$\varphi a\in pD_0$ and $v_p(n!)\le n-1$.
Approximation modulo $p$ extends the inclusion to $D$.
If $x=\varphi x$, choose $a\in A$ with $x\equiv a\pmod{pD}$.
Then $a^p-a\in pD$. It lies in $D_0$, so (LS.15) and
$\theta(D_0)\subset O$ imply
$(\theta a)^p-\theta a\in pO$. Its residue in the residue field
of $O$ therefore belongs to $\mathbb F_p$.
Choose $c\in\mathbb Z$ with $|\theta(a-c)|<1$.
Some power $\theta(a-c)^N$ is divisible by $p$ in $O$.
Lift its quotient by $\theta$, giving $b\in A$ with
$(a-c)^N-pb\in\ker\theta$. Its $p$-th power belongs to $pD_0$,
so $\bar a-c$ is nilpotent in $D/pD$. It also satisfies
$(\bar a-c)^p=\bar a-c$. A nilpotent satisfying this equation is
zero: multiply by the inverse of $1-(\bar a-c)^{p-1}$.
Hence $x\equiv c\pmod{pD}$.

As $p$ is regular and fixed by $\varphi$, the quotient $(x-c)/p$
is again fixed. Iterating with digits $0\le c_j<p$ gives
$x\equiv\sum_{j<n}p^jc_j\pmod{p^nD}$.
The digits converge to $z\in\mathbb Z_p$, and separation gives
$x=z$. Conversely Frobenius fixes $\mathbb Z_p$ by its Witt
definition (or by density of the integers and $p$-adic separation).

\paragraph{LS11: simultaneous theta vanishing and the Frobenius image.}
Two integral identities supply the fundamental intersection calculation:
\[
 \begin{split}
 &\theta(\varphi^na)=0\ (n\ge0),\ a\in A
                              \quad\Longrightarrow\quad \mu\mid a;\\
 &\varphi(D)\subset A+wD,
                 \qquad pw=\mu^{p-1}.
 \end{split}                                                   \tag{LS.30}
\]
For the first, let $\varepsilon_n=\varepsilon^{1/p^n}$,
$\xi_n=\sum_{i<p}[\varepsilon_{n+1}]^i$, and
$q_n=\prod_{j<n}\xi_j$. Then
$\mu=q_n([\varepsilon_n]-1)$.
Induction using $\ker\theta=(\xi)$ shows $q_n\mid a$:
after the first $n$ factors have been removed, apply $\theta\varphi^n$;
the removed factors have nonzero images $p$, so the next factor
must divide. Modulo $p$ this gives
\[
 |a_0|\le|\varepsilon-1|^{\,1-p^{-n}}\quad(n\ge0).
\]
Letting $n$ increase, the valuation-ring criterion implies
$(\varepsilon-1)\mid a_0$. Choose $b\in R$ with
$a_0=(\varepsilon-1)b$. Then $a=\mu[b]+pa'$.
Every $\theta\varphi^n(a')$ is zero, since the corresponding
values of $a$ and $\mu$ are zero and $p$ is regular in $O$.
Iteration and $p$-adic completeness prove $\mu\mid a$.

For the second identity, the geometric-sum polynomial gives
$\varphi\xi=pc+\mu^{p-1}=p(c+w)$, $c\in A$.
Thus
$\varphi(\xi^n/n!)=(p^n/n!)(c+w)^n\in A+wD$.
It follows for $D_0$ and then for $D$: successively write a dense
approximation to $y$, obtaining
$\varphi y=a_0+wz_0+p\varphi y_1$.
Iteration gives the convergent coefficients
$\sum p^ja_j\in A$ and $\sum p^jz_j\in D$, proving the claim.

\paragraph{LS12: integral Frobenius eigenvectors in the matching filtration.}
For $r\ge0$,
\[
 \varphi x=p^r x,\quad j(x)\in\operatorname{Fil}^r
 \quad\Longrightarrow\quad
 p^k x=z t^r\text{ for some }k\ge0, z\in\mathbb Z_p.
 \tag{LS.31}
\]
The case $r=0$ is LS10. For the induction step write
$\varphi x=p^{r+1}x=a+wz'$ by LS11. Every
$\theta_Dj(\varphi^n x)$ vanishes, as does
$\theta_Dj(\varphi^n w)$, the latter by
$p w=\mu^{p-1}$ and $\theta\varphi^n\mu=0$.
Thus $\theta\varphi^n a=0$ for every $n$, so $a=\mu b$.
With $t=\mu u$ as in LS7, set
\[
 x'=u^{-1}(pb+\mu^{p-2}z'),\qquad
                         p^{r+2}x=t x'.             \tag{LS.32}
\]
The exponent $p-2$ is nonnegative, also at $p=2$.
Apply Frobenius to the equality and cancel $pt$ in the domain $D$:
$\varphi x'=p^r x'$. Since $j(t)$ is a uniformizer and $p$ a
unit in $B_{\mathrm{dR}}^+$, (LS.32) gives
$j(x')\in\operatorname{Fil}^r$. Induction proves (LS.31), after
multiplication by one more power of $p$.

\paragraph{LS13: the actual crystalline localization and its fundamental intersection.}
Every generator $\mu_e=[e]-1$, with $e^\sharp=1$ and $e\ne1$,
is a unit in $D[1/t]$. Indeed LS7 gives
$t_e=\mu_eu_e\ne0$, $u_e\in D^\times$,
$\varphi t_e=pt_e$, and $j(t_e)\in\operatorname{Fil}^1$.
LS12 yields $p^kt_e=z t$ with $0\ne z\in\mathbb Z_p$.
Because $p$ is already a unit in $D[1/t]$, so is $z$; hence
$t_e$, and then $\mu_e$, is invertible. Conversely the source
localization inverts $\mu_\varepsilon$, and thus $t$.
The two localization maps are inverse and fix $D$. This identifies
the ring constructed here with the literal $B_{\mathrm{cris}}$
and the maps in LS8 with the source's maps. In particular the
constructed \texttt{Core} exists on the specified rings.

Now
\[
 B_{\mathrm{cris}}^{\varphi=1}\cap
        \iota^{-1}(B_{\mathrm{dR}}^+)=\mathbb Q_p.    \tag{LS.33}
\]
If $b$ belongs to the left side, write $bt^r=x\in D$.
Then $\varphi x=p^r x$ and $j(x)\in\operatorname{Fil}^r$.
Equation (LS.31) gives $p^kx=z t^r$, so $b=z/p^k$.
The reverse inclusion follows from the coefficient maps.
This proves exactly the \texttt{PeriodData.fund} field; a stronger
surjective fundamental exact sequence is not being asserted.

\paragraph{LS14: the de Rham invariants are exactly the finite base field.}
For every finite $K\subset\Omega$,
\[
                        B_{\mathrm{dR}}^{G_K}=o(K).       \tag{LS.34}
\]
First an invariant $z\in B_{\mathrm{dR}}^+$ with residue zero
vanishes. If $z=t^n y$, invariance and $g(t)=\chi(g)t$
give $g(y)=\chi(g)^{-n}y$. For $n>0$, its residue is a
nonzero Tate-twist invariant and is zero by LS3. Thus $t\mid y$.
Starting with $t\mid z$, induction makes $z$ divisible by every
power of $t$, and separation gives $z=0$.
More generally, if $g(y)=\chi(g)^n y$, $n\ge0$, then
\[
                y=t^n w,\qquad w\in(B_{\mathrm{dR}}^+)^{G_K}.
                                                               \tag{LS.35}
\]
For $n>0$, LS3 kills the residue of $y$; divide by $t$ and
induct with exponent $n-1$. The case $n=0$ is tautological.

Now clear the denominator of an invariant $b\in B_{\mathrm{dR}}$:
$y=t^n b\in B_{\mathrm{dR}}^+$ for some $n\ge0$.
Equation (LS.35) gives $y=t^n w$ with $w$ invariant, so $b=w$.
By LS1 its residue belongs to $K$. Subtract its unique algebraic
lift $o(k)$ from LS5. The difference is invariant with zero
residue, hence is zero by the first paragraph. Conversely $o(k)$
is invariant for $k\in K$, by equivariance of the algebraic
section. This proves (LS.34), including the exclusion of negative
filtration invariant elements.

\paragraph{LS15: the crystalline invariants are exactly the unramified base field.}
For the same $K$,
\[
                         B_{\mathrm{cris}}^{G_K}=u(K_0).
                                                               \tag{LS.36}
\]
Let $b$ be invariant. Frobenius commutes with $G_K$, and LS14
therefore gives elements $y_i\in K$ with
$\iota(\varphi^i b)=o(y_i)$. If $q\in\mathbb Q_p[X]$
annihilates $y_0$, injectivity of $\iota$ gives $q(b)=0$;
Frobenius fixes $\mathbb Q_p$, so $q(\varphi^i b)=0$ for
all $i$. The nonzero minimal polynomial of $y_0$ has finitely
many roots. Thus two iterates of $b$ agree, and injectivity of
Frobenius (LS9) permits cancellation of their common initial
iterate. There is an integer $N>0$ with $\varphi^N b=b$.

Choose a primitive root $\zeta\in\mathbb Q_p^{\mathrm{ur}}$
of order $p^N-1$. Its Frobenius orbit
$\zeta,\sigma\zeta,\ldots,\sigma^{N-1}\zeta$ consists of
$N$ distinct points: $p^i-1<p^N-1$ for $0<i<N$.
When $N=1$, this assertion has no distinctness condition, including
$p=2$, where $\zeta=1$. For every integer $n\ge0$, set
\[
 S_n=\sum_{i=0}^{N-1}u((\sigma^i\zeta)^n)\,\varphi^i b.
                                                               \tag{LS.37}
\]
Frobenius permutes the summands. Their images under $\iota$
are algebraic coefficients in $B_{\mathrm{dR}}^+$.
The fundamental intersection LS13 therefore gives
$S_n\in\mathbb Q_p$. Consequently every polynomial-weighted
sum in (LS.37), with coefficients in
$\mathbb Q_p^{\mathrm{ur}}$, belongs to $u(\mathbb Q_p^{\mathrm{ur}})$.
Take
$P(X)=\prod_{j=1}^{N-1}(X-\sigma^j\zeta)$.
All terms except the one with $i=0$ vanish, so
$u(P(\zeta))b\in u(\mathbb Q_p^{\mathrm{ur}})$.
Its nonzero coefficient is invertible, whence $b=u(c)$ for some
$c\in\mathbb Q_p^{\mathrm{ur}}$. Comparing with
$\iota b=o(y_0)$ and using injectivity of $o$ gives $c=y_0\in K$.
Thus $c\in K_0$. Conversely $u(K_0)$ is fixed by $G_K$ by
LS8. This proves the exact invariant-field assertion without
assuming a crystalline comparison theorem for representations.

\paragraph{LS16: all fields of the actual finite-extension period datum.}
For every finite $K\subset\Omega$, the preceding construction
gives the literal datum
\[
 \operatorname{PeriodData}
   (G_K,\mathbb Q_p,K_0,K,B_{\mathrm{cris}},B_{\mathrm{dR}})
                                                               \tag{LS.38}
\]
on the same two period rings, with no new choice of a model for
each $K$. We spell out the full record verification. Both $K$
and $K_0$ are finite-dimensional over $\mathbb Q_p$; the latter
is a subspace of $K$. The residue Frobenius $\sigma$ of LS8
restricts to a bijective $\mathbb Q_p$-map on $K_0$.
The restricted $G_K$-actions on $B_{\mathrm{cris}}$ and
$B_{\mathrm{dR}}$ are respectively $K_0$- and $K$-linear,
because those constants are fixed. Their inverses are the actions
of $g^{-1}$, and the group laws hold on coordinates and hence
on completions and localizations. Their invariant rings are
exactly the coefficient fields by LS14--LS15.

Frobenius is an injective $\mathbb Q_p$-algebra map by LS9;
it acts as $\sigma$ on $K_0$ and commutes with the group action
by LS8 and the commuting Witt operations. The map $\iota$ is an
injective $K_0$-algebra map and is equivariant. For each integer
$a$, take the actual $\operatorname{Fil}^a=\xi^aB_{\mathrm{dR}}^+$.
This is a $K$-subspace: $o(K)\subset B_{\mathrm{dR}}^+$ and
multiplication preserves the indicated power. It contains $1$
in degree zero, is decreasing, exhaustive, separated,
multiplicative and $G_K$-stable by LS5.
Take the restriction of $\chi$, valued in $\mathbb Q_p^\times$,
and the same $t$ of LS7. It is a unit in $B_{\mathrm{cris}}$,
$g(t)=\chi(g)t$, $\varphi(t)=pt$, and its image lies in
$\operatorname{Fil}^1\setminus\operatorname{Fil}^2$ because
it is a uniformizer. Finally LS13 gives the precise fundamental
intersection field. These are all fields of the source record.

The equalities in \texttt{IsRealPD} are literal identifications:
the prime is $p$; the two actions, $\varphi$, $\iota$, $t$,
$\chi$ and every filtration subset are those just specified;
and the restricted $\sigma$ is the same unramified Frobenius.
In particular, changing $K$ changes the coefficient subfields and
restricted group, but does not change the filtration subsets in
$B_{\mathrm{dR}}$.

\paragraph{LS17: the complete local prime, Ainf tie, and all-prime stage.}
Fix one construction of LS1--LS16 for each rational prime $p$.
It yields simultaneously
\[
 \exists(P:\texttt{LocPrime}\ p)\ (AI:\texttt{LocAinf}\ p\ P),
                 \texttt{LocTied}\ p\ P\ AI,\qquad
                 \texttt{LocStage}.                         \tag{LS.39}
\]
Here is the remaining field audit, including the ties to actual
rings. Take the underlying rings of $P$ to be the actual
$B_{\mathrm{cris}}$ and $B_{\mathrm{dR}}$, with their
$\mathbb Q_p$, unramified, and algebraic coefficient structures
from $u$ and $o$. If the source presents these as renamed types,
transport every operation along the identity ring equivalences
$e_B,e_D$. The datum for every finite $K$ is LS16.
The global actions, Frobenius, inclusion, $t$, and $\chi$ are
those of LS8; therefore every restriction equality in
\texttt{LocPrime} holds. Its equality of filtration subsets for
any two finite $K,K'$ follows from LS16. The actions on
unramified and algebraic constants are the ones in LS8 and LS5.
Frobenius on the valuation ring of $K_0$ reduces to $x\mapsto x^p$.

The cyclotomic clause includes its numerical approximation:
for $g\in G$ and $m\ge0$, choose the least nonnegative integer
$a$ representing $\chi(g)\bmod p^m$. Then
\[
 |\chi(g)-a|\le p^{-m},\qquad g(\zeta)=\zeta^a
                   \quad(\zeta^{p^m}=1).                  \tag{LS.40}
\]
The compatible primitive system $\zeta_m$ is the one chosen in
LS2. For every finite extension $L/\mathbb Q_p$, use its unique
extended absolute value and put $v(x)=-\log_p|x|$ for $x\ne0$,
$v(0)=\infty$. Algebraicity makes the nonzero values rational,
and $v(p)=1$. A finite complete DVR extension, as in NC1--NC2,
has $v(x)\ge0$ exactly on the integral closure of $\mathbb Z_p$.
This verifies both normalized valuation fields of $P$.

For $AI$, extend $G$ isometrically from $\Omega$ to $C$ and
act coordinatewise on $R=\varprojlim O/p$. The coordinate
formula holds because each action preserves $pO$ and commutes
with Frobenius. Take the same $\varepsilon$, its exact sharp-root
system, and the actual map $A\hookrightarrow B_{\mathrm{cris}}$.
This map is injective by LS5--LS6, is equivariant, and commutes
with Witt Frobenius by its construction. Its image under
$\iota$ is in filtration zero. For every finite $K$,
\[
 \iota(a)\in\operatorname{Fil}^1
      \quad\Longleftrightarrow\quad\theta(a)=0
      \quad\Longleftrightarrow\quad\xi\mid a\text{ in }A.
                                                               \tag{LS.41}
\]
This is the kernel calculation LS4, not merely divisibility in
a period field. For every prime-to-$p$ root $c$ and every tilt
representative satisfying the specified root-coordinate condition,
the Teichm\"uller image is $u(c)$ by LS8: uniqueness of the
root in the complete DVR proves the equality for all such choices.
These verify all \texttt{LocAinf} fields.

The tie uses $e_B,e_D$ just specified. It sends the prescribed
$\mathbb Z_p$ constants to the actual constants because $o$
and $u$ agree on them. Each algebraic $x\in\Omega$ has the
lift $o(x)$ with the required residue. Actions and Frobenius
are their actual coordinate constructions. The transported
inclusion satisfies \texttt{IsCrisDR}: its restriction to $D$
is the actual $j$, with the bounds proved in LS6, and its
extension is the localization map LS8--LS13. Filtrations agree
for every $K$ and degree. The same $t$ maps to $j(t)$, whose
formal logarithm identity was proved in LS7. Finally the
$A_{\inf}$ inclusion is the actual one. Thus every field of
\texttt{LocTied}, including its integral and logarithmic clauses,
has been verified. Choosing these triples for all primes gives
\texttt{LocStage}; this uses the constructions at every prime,
even when a later application has an empty local index type.

\paragraph{LS18: the global cyclotomic, residue, and prime-label suppliers.}
For the remainder fix the source $B:\texttt{Base}$, its host $H$,
$C:\texttt{Cols}(B)$, and a row object $W:\texttt{RowsRest}(B,C)$.
Write $r=B.R.r\ge2$ and $f=B.R.f>0$.
The nonempty pieces of the bare \texttt{NData} and \texttt{NTypes}
construction have the following elementary suppliers.

For each rational prime $p$, let $\Omega_H=\overline{H.E}$ and
$S_p$ be the finite set of primes of $H.E$ above $p$.
Finiteness follows from the finite factorization of the nonzero
ideal $(p)$. Choose a valuation above each such prime.
The cyclotomic character on
$\operatorname{Gal}(\Omega_H/H.E)$ is unramified outside $S_p$.
Indeed inertia at residue characteristic $q\ne p$ fixes the
residue of each $p^n$-th root of unity. The quotient of two such
roots with the same residue is $1$: if $z^m=1$, $q\nmid m$,
and $z\equiv1$, cancellation of $z-1$ in the geometric sum
would force its residue $m$ to be zero unless $z=1$.
Thus inertia fixes every such root, and the inverse-limit
character is trivial on inertia. It descends through the quotient
unramified outside $S_p$. Together with the chosen valuations,
this is the actual \texttt{LData} used in the source. Set
\[
 s_H(p,e)=
 \begin{cases}
  (\operatorname{shaLen}(L_p,e).\mathrm{toNat}:\mathbb Q),
                                      &p\text{ prime},\\
  0,&p\text{ not prime}.
 \end{cases}                                                \tag{LS.42}
\]
If the extended-natural length is the finite integer $n$,
then $s_H(p,e)=n$. This proves the exact conditional
\texttt{sha\_eq} field. It neither proves finiteness nor rules
out infinite length, whose \texttt{toNat} is zero.

There is a complex-conjugation lift on $\Omega_H$. Embed this
algebraic closure into $\mathbb C$ over $\mathbb Q$.
Its image is the field of algebraic numbers, stable under complex
conjugation. Conjugate through the embedding; injectivity gives
an involution, and the host's CM embedding identity identifies
its restriction to $H.E$ with the prescribed conjugation.

For each rational prime $g$, there is a multiplicative map
$\overline{\mathbb F}_g\to\mathbb C$ injective on nonzero
elements. To construct it directly, exhaust the algebraic closure
by $\mathbb F_{g^{n!}}$. Their multiplicative groups are cyclic
and nested. Embed each into the complex roots of unity, extending
the previous embedding: a prescribed primitive $d$-th root
extends to a primitive $N$-th root whenever $d\mid N$, by the
surjectivity of
$(\mathbb Z/N)^\times\to(\mathbb Z/d)^\times$, proved prime
by prime and by the Chinese remainder theorem. Put the value at
zero equal to zero. This supplies the exact multiplicative
\texttt{teichC} interface; additivity is not required.
For any prime $\ell$, choose a reduction from the ring of all
algebraic integers in $\mathbb C$ to
$\overline{\mathbb F}_\ell$. Such a map exists by taking a
maximal ideal above $(\ell)$ in this integral extension of
$\mathbb Z$ and embedding its algebraic residue field in the
algebraic closure. Every nonzero value of the preceding map is
a root of unity, hence integral. Composing gives the required
multiplicative map between residue algebraic closures. It is
injective on elements of order prime to $\ell$: their quotient
reducing to $1$ is $1$ by the same geometric-sum argument.
This proves exactly the restricted injectivity in
\texttt{iota\_inj}, without asserting an embedding of fields of
different characteristics.

Finally the source congruences $B.g\_cond$ make each $g_k$
split completely in $H.E$, by the host splitting calculation
already reconstructed in AH. The Galois group acts transitively
on primes above $g_k$; equality of their number with its order
makes this action free. Host conjugation is a nonidentity
involution, so its action on these primes has no fixed point.
There are $2f$ primes. Choose one prime from each conjugate pair,
giving an injective map $\operatorname{Fin}(f)$ into the primes
above $g_k$, with no two selected primes conjugate. The quadratic
field $H.K$ has discriminant $-H.d$; its identity extension over
$H.K$ is abelian, its ramified rational primes divide the source
conductor, and the same congruence makes $g_k$ split there.
These verify the selected-prime and ray-field clauses without
requiring a nontrivial class-field extension.

\paragraph{LS19: an actual tame inertia generator at every selected place.}
Let $V$ be a valuation of $\Omega_H$ above a finite place of the
number field $H.E$, and $I$ its inertia subgroup. There is
$\gamma\in I$ satisfying the source's \texttt{TameGenAt}:
\[
 I\subset\overline{\langle\gamma,P_V\rangle},\qquad
 P_V=\{\tau\in I:|\tau(y)/y-1|<1\text{ for every }y\ne0\}.
                                                               \tag{LS.43}
\]
Here $P_V$ is the explicitly defined wild subgroup, so (LS.43)
does not presuppose a tame-inertia classification theorem.

We use a compactness lemma. Suppose a compact group $I$ has
continuous maps $\theta_n$ to finite cyclic groups, with nested
kernels. The set of elements whose $\theta_n$-image generates
$\theta_n(I)$ is nonempty and closed. These sets are nested:
a generator maps to a generator in every quotient.
Compactness gives an element $\gamma$ in their intersection.
For each $\sigma\in I$, the sets
\[
 \overline{\langle\gamma\rangle}\cap
                       \theta_n^{-1}(\theta_n\sigma)
\]
are again nonempty, closed and nested, because a power of
$\gamma$ has any prescribed image in the finite cyclic group.
Their intersection contains an $h$ with
$h^{-1}\sigma\in\bigcap_n\ker\theta_n$.
Thus it suffices to identify this common kernel with $P_V$.

The restriction of the valuation to $\mathbb Q$ is controlled
by a single prime $q$: all other rational primes have norm $1$.
For any algebraic $y\ne0$, its polynomial equation has at least
two terms of largest norm; otherwise their sum could not be zero.
Equating those two norms gives
$|y|^a=|q|^b$ for some $a>0$, $b\in\mathbb Z$.
Choose $r_n\in\Omega_H$ with $r_n^{n!}=q$.
For $\sigma\in I$, take the residue class of
$\sigma(r_n)/r_n$ in the group of norm-one elements modulo
principal units, and call it $\theta_n(\sigma)$.
It is a homomorphism, because inertia acts trivially on residues.
Its image is contained in the reduction of the finite group of
$n!$-th roots of unity, hence is finite cyclic. Its fibers are
closed: the action on the fixed algebraic element $r_n$ factors
through a finite Galois extension. The kernels are nested.
Indeed $r_{n+1}^{n+1}/r_n$ has norm $1$, and the action of
inertia on this quotient changes it only by a principal unit.

If $\tau$ belongs to every kernel and $y\ne0$, choose $a,b$
as above and $n$ with $a\mid n!$. The element
$r_n^{b n!/a}$ has the same norm as $y$, and its inertia
coboundary is a principal unit. The quotient of $y$ by this
element is a unit, whose inertia coboundary also has residue $1$.
Thus $|\tau(y)/y-1|<1$, so $\tau\in P_V$.
The reverse inclusion follows by testing $y=r_n$.
Inertia is a closed subgroup of the profinite absolute Galois
group and is compact. The lemma therefore proves (LS.43).
Applying it to every selected prime from LS18 supplies the
literal \texttt{tg}, \texttt{tg\_mem}, and \texttt{tg\_gen}
fields; no selected-place generator is omitted by an empty-index
argument.

\paragraph{LS20: the full empty-index algebraic and automorphic records.}
Suppose
\[
 W.\mathrm{shaLen}=s_H,\qquad
 W.\eta_J(j,w)=0\ (1\le j\le m),\qquad
 W.e\mathrm{Op}=0,\quad W.eM_t=0.                    \tag{LS.44}
\]
The source constructs the following records on exactly these
$B,C,W$. Take $\texttt{NData.Om}=\Omega_H$, the chosen
valuations and $L_p$ from LS18, and put
\[
 \mathrm{Loc}=\varnothing,\quad\mathrm{locs}(j,w)=\varnothing,
 \quad\mathrm{PlG}(k)=\operatorname{Fin}(f),\quad
 \mathrm{lab}=(X-1)^r,\quad\chi_T=1.                  \tag{LS.45}
\]
The split-place and cardinality fields come from LS18 and the
\texttt{RowsRest} finite-group cardinal formula. The polynomial
is monic of degree $r$ with nonzero constant term. Its Sha clause
is precisely the finite-value assertion in LS18. Every field
indexed by $i\in\mathrm{Loc}$, including its local DVR,
representation, eigenvector, auxiliary-prime, and local-generator
clauses, is a function on the empty type. This supplies all of
\texttt{NData}, without introducing a local representation.

For clarity the nonempty part of \texttt{NAut} is also explicit.
At each place let a local parameter be a finite multiset of
nonzero complex numbers. Its dimension is cardinality; zero
is the empty multiset; addition is multiset sum. Twisting by a
real $s$ multiplies every entry by $\mathrm{Norm}(w)^s$;
character twisting multiplies by the product of the entries of
the character parameter. These scalars are nonzero. Set
\texttt{unr} and \texttt{mono0} identically true,
\texttt{frobEig} equal to that multiset, and temperedness equal
to the condition that all entries have absolute value $1$.
Set the inertia polynomial and \texttt{items2} equal to
$(X-1)^{\mathrm{dimension}}$. Multiset addition and scaling
verify the dimension, sum, twist, Frobenius, and inertia laws;
$(X-1)^{a+b}=(X-1)^a(X-1)^b$ gives their polynomial identity.

A global parameter of rank $a$ is an arbitrary family of these
local multisets of size $a$, together with one complex multiset
of size $a$ for each embedding of $H.E$ into $\mathbb C$.
Conjugate self-duality means inversion of the local multiset at
the conjugate place. Regular algebraicity means that each
infinite multiset has no repetitions and consists of elements
$n+(a-1)/2$, $n\in\mathbb Z$. The global twist scales the
finite entries and adds the sum of the twisting infinite entries.
Inversion of a product proves preservation of self-duality;
cardinality and the one-element case give the remaining stated
twist laws. Choose both \texttt{etaT} and its inverse to have
finite multiset $\{-1\}$ and infinite multiset $\{0\}$.
The inert-place and inverse clauses hold, and at a split pair
$(-1)(-1)=1$. Use the actual selected primes of LS18 for
\texttt{plGw}. All remaining \texttt{NAut} clauses have an
index in $\mathrm{Loc}$ and are vacuous. These are literal record
models; they are not asserted to classify automorphic representations.

In \texttt{NBC} take both representation and weight indexing
types empty. In \texttt{NTw} take the actual CM-conjugation
lift from LS18 and $\mathrm{Tw}=\varnothing$. Every other
field is indexed by $\mathrm{Tw}$ or $\mathrm{Loc}$, including
coefficient-DVR, character, Galois, and local-global ties.
For \texttt{NExp} use the zero distinguished subspace and zero
coefficient function, and use the empty functions for its twist
clauses. All \texttt{NTh} fields are twist-indexed and are
supplied by the same empty functions. This verifies the actual
constructors \texttt{nBC}, \texttt{nTw}, \texttt{nExp}, and
\texttt{nTh} with their exact common \texttt{NData}; no
existence assertion over a nonempty twist set is used.

\paragraph{LS21: every type, projector, bridge, and row tie in the bare NStage.}
Take the model of LS20 and an arbitrary supplied
$L:\texttt{LocStage}$. In \texttt{NTypes}, set the ray field
to $H.K$ with its identity extension, and use the conductor,
splitting, and disjoint \texttt{BB} data of $B$, checked in LS18.
Use the block multiset $\{(r,1)\}$ and the trivial residual
representation. Its kernel is the whole Galois group and is open.
Use the actual tame generators of LS19 and the multiplicative
residue maps of LS18. The inertia characteristic polynomial is
correct because the roots of $(X-1)^a$ are $a$ copies of $1$,
and every multiplicative map sends $1$ to $1$.

Define, in the relevant finite group algebra,
\[
 e_f=\begin{cases}1,&f=(X-1)^r,\\0,&\text{otherwise}.
       \end{cases}                                      \tag{LS.46}
\]
These elements are central, idempotent, and pairwise orthogonal.
The finite set of monic degree-$r$ polynomials with nonzero
constant coefficient contains $(X-1)^r$, and summing the $e_f$
over it gives $1$. Central actions on the ambient and middle
spaces are the augmentation map followed by scalar identity
endomorphisms; this is an algebra homomorphism. Take the type
index to be a singleton, every type label to be $(X-1)^r$,
each local type multiset to consist of $r$ copies of $1$, and
both type subspaces to be the whole space. These choices verify
the remaining nonempty \texttt{NTypes} clauses. Its two ties
quantified by a local index are vacuous.

A good polynomial is squarefree. Since $r\ge2$, it cannot be
$(X-1)^r$, so its $e_f$ is zero. There is a selected prime
because $f>0$. Each product indexed by a good tuple therefore
has a zero factor, and the sum $e_0$ over good tuples is zero.
This is the full calculation underlying SE13. By (LS.44), the
augmentation image of $e_0$ equals the required zero operators
on both ambient and middle spaces. Every coefficient of $e_0$
is zero, so the required coefficient-field equality with
\texttt{NExp.eLL0} also holds.

Put all these records together and retain $L$ literally as the
local-stage field. The bridge fields \texttt{nbr},
\texttt{nbr\_ind}, and \texttt{self\_good} have local indices;
\texttt{ntwbr} and \texttt{ntwbr\_ind} have twist indices.
Their types are empty. The local, base-change, residual,
middle-parameter, $Z$, and $Z\mathrm{good}$ ties likewise
have local indices. Finally, for $1\le j\le m$,
\[
 W.\eta_J(j,w)=0
       =\sum_{i\in\mathrm{locs}(j,w)}\eta_{\mathrm{Loc}}(i).
                                                               \tag{LS.47}
\]
Thus no field remains unchecked, and we obtain the literal endpoint
\[
 \exists N:\texttt{NStage}(B,C,W),\qquad N.\mathrm{toLocStage}=L.
                                                               \tag{LS.48}
\]
In particular LS17 supplies such an $L$. For arbitrary $W_0$,
SE13's exact update $W=\mathrm{adapt}(W_0)$ satisfies (LS.44).
This proves \texttt{exists\_nStage\_of\_rows},
\texttt{exists\_nStage\_adapt}, and their existential row-object
version, including preservation of the same prescribed $L$.

\paragraph{LS22: the bare Later record now exists, with its exact boundary.}
Let $I$ be the supplied common column input, let $l:I.\mathrm{Laws}$,
and suppose \texttt{HasLetters} holds for the same $B$.
Set $C=I.\mathrm{cols}(l)$, and take exactly the row model
$W_0$ and update $W=\mathrm{adapt}(W_0)$ of SE14.
By LS17 and LS21, there is an
$N:\texttt{NStage}(B,C,W)$. SE6--SE12 supply the other
three auxiliary stages on this same $B,C,W$, and SE14 now gives
\[
                         \operatorname{Nonempty}
                              (\texttt{Later}(B,C)).       \tag{LS.49}
\]
This discharges the full bare \texttt{LocStage}/\texttt{NStage}
existence input left open there, including the specified integral
period rings and their ties. It neither replaces an already fixed
later stage by a different one nor supplies extra properties that
are absent from the bare record.

The distinction from arithmetic choice is forced by the exact
index types. The constructed stage has both
$\mathrm{Loc}=\varnothing$ and $\mathrm{Tw}=\varnothing$.
The arithmetic \texttt{NChoice} needs a function from every
coefficient-field prime to $\mathrm{Tw}$, even at a reference
weight. A coefficient number field has a prime above $2$, so
this domain is nonempty and such a function cannot exist for
this particular stage. The empty-local-index hypotheses used in
RC do not remove that requirement: its additional twist-reading
hypothesis must still be supplied. No crystalline, de Rham, or
Wach comparison for a nonempty family of geometric representations
has been deduced merely from the period-ring construction.
Similarly the model's zero projectors do not furnish positive rank
bounds or the ties \texttt{RankTie}, \texttt{MuIsW}, and the
same-object arithmetic data needed by later quantitative arguments.
These are the precise continuing boundaries after (LS.49).
